\section{Кросс-энтропия. Дивергенция Кульбака-Лейблера}
\localtableofcontents

\subsection{Кросс-энтропия}

Перед тем как вводить дивергенцию Кульбака–Лейблера, полезно рассмотреть более простую величину --- кросс-энтропию. Она измеряет \enquote{среднюю удивлённость} при использовании модели $q$ для кодирования данных, порождённых истинным распределением $p$.

\begin{definition}[Кросс-энтропия для дискретных распределений]
	Пусть $p$ и $q$ --- два распределения на одном и том же конечном множестве $\mathcal{X}$ (вероятностные функции). Кросс-энтропией из $p$ в $q$ называется величина
	\[
		\CE_m(p | q) = -\sum_{x \in \mathcal{X}} p(x) \log_m q(x).
	\]
\end{definition}

\begin{remark}[Связь с энтропией Шеннона]
	Если $q = p$, то кросс-энтропия превращается в обычную энтропию Шеннона:
	\[
		\CE_m(p | p) = -\sum_x p(x) \log_m p(x) = \H_m(p).
	\]
\end{remark}

\begin{proposition}[Нижняя граница кросс-энтропии]
	\label{prop:cross-entropy-min}
	Для любых распределений $p$ и $q$ на $\mathcal{X}$ выполнено:
	\[
		\CE_m(p | q) \geq \H_m(p),
	\]
	причём равенство достигается тогда и только тогда, когда $p = q$.
\end{proposition}
\begin{proof}
	Переформулировка неравенства Гиббса (лемма \ref{lm:gibbs-ineq}) 
\end{proof}

\begin{definition}[Дифференциальная кросс-энтропия]
	Пусть $p$ и $q$ --- две плотности распределений на $\mathbb{R}^d$, причём носитель $p$ содержится в носителе $q$ (т.е. $q(\vx) = 0 \Rightarrow p(\vx) = 0$ для почти всех $\vx$). Дифференциальной кросс-энтропией из $p$ в $q$ называется величина
	\[
		\CE_m(p | q) = -\int_{\mathbb{R}^d} p(\vx) \log_m q(\vx) \, d\vx.
	\]
\end{definition}
\begin{proposition}[Неравенство Гиббса для дифференциальной кросс-энтропии]
	\label{prop:diff-cross-entropy-gibbs}
	Для любых плотностей $p$ и $q$ выполнено, таких что $\CE_m(p | q)$ определена:
	\[
		\CE_m(p | q) \geq h_m(p),
	\]
	где $h_m(p)$ --- дифференциальная энтропия распределения $p$. Равенство достигается тогда и только тогда, когда $p = q$ почти всюду.
\end{proposition}
\begin{proof}
	Следует из неравенства Гиббса (лемма \ref{lem:gibbs-continuous}). Условие равенства также следует из леммы.
\end{proof}

\subsection{Дивергенция Кульбака-Лейблера}

\begin{definition}[Дивергенция Кульбака–Лейблера]
	Пусть $P$ и $Q$ --- два распределения на одном и том же измеримом пространстве, причем $P$ абсолютно непрерывно относительно $Q$ ($P \ll Q$). Дивергенция Кульбака–Лейблера (или относительная энтропия) определяется как:
	\[
		\KL(P | Q) = \int \log_m(\od{P}{Q}) \dif P,
	\]
	где $\od{P}{Q}$ --- плотность $P$ относительно $Q$ (существует по теореме Радона-Никодима). \\
	В случае дискретных распределений с вероятностными функциями $p(x)$ и $q(x)$:
	\[
		\KL(p | q) = \sum_x p(x) \log_m \frac{p(x)}{q(x)}.
	\]
	В случае непрерывных распределений с плотностями $p(\vx)$ и $q(\vx)$:
	\[
		\KL(p | q) = \int p(\vx) \log_m \frac{p(\vx)}{q(\vx)} \dif \vx.
	\]
\end{definition}

\begin{proposition}[Неотрицательность дивергенции Кульбака–Лейблера]
	Для любых распределений $P$ и $Q$ таких, что $P \ll Q$, выполнено:
	\[
		\KL(P | Q) \geq 0,
	\]
	причём равенство достигается тогда и только тогда, когда $P = Q$ почти всюду.
\end{proposition}
\begin{proof}
	В дискретном случае утверждение следует непосредственно из неравенства Гиббса (лемма \ref{lem:gibbs}). В непрерывном случае — из леммы \ref{lem:gibbs-continuous}, применённой к плотностям $p$ и $q$. В общем случае мер произвольной природы доказательство использует неравенство Йенсена для вогнутой функции $\log_m$ и тот факт, что $\int \frac{dP}{dQ} \, dQ = 1$.
\end{proof}

\begin{proposition}[Связь с энтропией]
	Для дискретного распределения $p$ на конечном множестве $\mathcal{X}$ мощности $|\mathcal{X}|$ выполнено:
	\[
		\H_m(p) = \log_m |\mathcal{X}| - \KL(p | u),
	\]
	где $u$ --- равномерное распределение на $\mathcal{X}$.
\end{proposition}
\begin{proof}
	Вычислим:
	\[
		\KL(p | u) = \sum_x p(x) \log_m \frac{p(x)}{1/|\mathcal{X}|} = \sum_x p(x) \log_m p(x) + \log_m |\mathcal{X}| \sum_x p(x) = -\H_m(p) + \log_m |\mathcal{X}|.
	\]
	Откуда непосредственно получаем требуемое равенство.
\end{proof}

\begin{remark}
	Дивергенция Кульбака–Лейблера не является метрикой, так как она не симметрична ($\KL(P\|Q) \neq \KL(Q\|P)$ вообще говоря) и не удовлетворяет неравенству треугольника. Тем не менее, она обладает важным свойством неотрицательности и играет ключевую роль в теории информации, статистике и машинном обучении как мера «расстояния» между распределениями.
\end{remark}

\subsection{Взаимная информация между случайными величинами}

\begin{definition}[Взаимная информация]
	Для случайных величин $\xi$ и $\eta$ с совместным распределением $\P_{\xi, \eta}$ и маргинальными $\P_{\xi}$, $\P_{\eta}$ взаимная информация определяется как:
	\[
		\I_m(\xi;\eta) = \KL(\P_{\xi; \eta} | \P_\xi \otimes \P_{\eta}).
	\]
	В дискретном случае:
	\[
		\I_m(\xi;\eta) = \sum_{x,y} p_{\xi, \eta}(x,y) \log_m \frac{p_{\xi, \eta}(x,y)}{p_{\xi}(x)p_{\eta}(y)}.
	\]
	В непрерывном случае:
	\[
		\I_m(\xi;\eta) = \iint f_{\xi; \eta}(\vx,\vy) \log_m \frac{f_{\xi, \eta}(\vx,\vy)}{f_{\xi}(\vx)f_{\eta}(\vy)} \dif\vx\dif\vy.
	\]
\end{definition}

\begin{proposition}[Эквивалентные формы]
	Взаимная информация может быть выражена через энтропии:
	\[
		\I(X;Y) = \H(X) - \H(X|Y) = \H(Y) - \H(Y|X) = \H(X) + \H(Y) - \H(X,Y).
	\]
	
	Для непрерывных случайных величин аналогично:
	\[
		\I(X;Y) = h(X) - h(X|Y) = h(Y) - h(Y|X) = h(X) + h(Y) - h(X,Y).
	\]
\end{proposition}

\begin{proof}
	Непосредственно следует из определений энтропии и условной энтропии.
\end{proof}


